\vspace{0.5em}\noindent\textbf{Workshop}\par

\vspace{0.5em}\noindent\textbf{Overview}\par

This workshop presents the complete deployment process of the Internship
Application Tracker on AWS.

The workshop covers source code preparation, cloud infrastructure
provisioning, database deployment, application deployment, AI
processing, event-driven notification, monitoring, testing,
troubleshooting, and resource clean-up.

The final architecture uses:

\begin{itemize}
\tightlist
\item
  Amazon CloudFront and Amazon S3 for frontend delivery.
\item
  Application Load Balancer for public API and realtime traffic.
\item
  Amazon EKS for backend, chat, worker and AI adapter workloads.
\item
  Amazon RDS PostgreSQL for transactional business data.
\item
  Amazon ElastiCache Redis for realtime Socket.IO pub/sub.
\item
  Amazon DynamoDB for chat persistence and Lambda event deduplication.
\item
  Amazon SQS for asynchronous event delivery.
\item
  AWS Lambda and Amazon SES for notifications.
\item
  Amazon SageMaker for AI inference.
\item
  GitHub Actions and AWS OIDC for continuous deployment.
\end{itemize}

\vspace{0.5em}\noindent\textbf{Learning objectives}\par

After completing this workshop, the reader should be able to:

\begin{enumerate}
\def\labelenumi{\arabic{enumi}.}
\tightlist
\item
  Prepare a multi-service application for AWS deployment.
\item
  Deploy containerized services to Amazon EKS.
\item
  Deploy a static React frontend through Amazon S3 and CloudFront.
\item
  Connect EKS workloads to RDS, Redis, DynamoDB, SQS and SageMaker.
\item
  Implement asynchronous processing with workers and transactional
  outbox.
\item
  Process SQS events with AWS Lambda.
\item
  Implement idempotent Lambda processing using DynamoDB.
\item
  Configure CI/CD using GitHub Actions and AWS OIDC.
\item
  Validate the complete system through end-to-end testing.
\item
  Diagnose common networking, IAM and deployment failures.
\end{enumerate}

\vspace{0.5em}\noindent\textbf{Target audience}\par

This workshop is intended for:

\begin{itemize}
\tightlist
\item
  Cloud engineering interns.
\item
  Backend and DevOps students.
\item
  Developers learning Amazon EKS.
\item
  Developers implementing event-driven AWS architectures.
\item
  Students preparing an AWS internship or graduation project report.
\end{itemize}

\vspace{0.5em}\noindent\textbf{Architecture summary}\par

The browser connects to Amazon CloudFront through HTTPS.

CloudFront routes:

\begin{itemize}
\tightlist
\item
  static frontend requests to Amazon S3
\item
  \texttt{/\allowbreak{}api/\allowbreak{}*} requests to the backend through the Application Load
  Balancer
\item
  \texttt{/\allowbreak{}chat/\allowbreak{}*} requests to the chat service
\item
  \texttt{/\allowbreak{}socket.\allowbreak{}io/\allowbreak{}*} requests to the realtime Socket.IO service
\end{itemize}

The backend, chat service, processing worker, AI service and outbox
dispatcher run inside Amazon EKS.

PostgreSQL stores transactional business data and processing jobs.
DynamoDB stores chat data. Redis provides pub/sub between chat pods.

The processing worker calls the AI service, which invokes the SageMaker
endpoint.

The outbox dispatcher publishes business events to Amazon SQS. AWS
Lambda consumes these events, performs deduplication in DynamoDB,
archives events to Amazon S3 and sends notifications through Amazon SES.

\vspace{0.5em}\noindent\textbf{Chapters}\par

{
\begingroup
\small
\setlength{\tabcolsep}{3pt}
\renewcommand{\arraystretch}{1.15}
\begin{longtable}{@{}
  >{\raggedright\arraybackslash}p{(\linewidth - 6\tabcolsep) * \real{0.2500}}
  >{\raggedright\arraybackslash}p{(\linewidth - 6\tabcolsep) * \real{0.2500}}
  >{\raggedright\arraybackslash}p{(\linewidth - 6\tabcolsep) * \real{0.2500}}
  >{\raggedright\arraybackslash}p{(\linewidth - 6\tabcolsep) * \real{0.2500}}@{}}
\toprule\noalign{}
\begin{minipage}[b]{\linewidth}\raggedright
Chapter
\end{minipage} & \begin{minipage}[b]{\linewidth}\raggedright
Title
\end{minipage} & \begin{minipage}[b]{\linewidth}\raggedright
Description
\end{minipage} & \begin{minipage}[b]{\linewidth}\raggedright
Status
\end{minipage} \\
\midrule\noalign{}
\endhead
\bottomrule\noalign{}
\endlastfoot
5.1 & Overview & Workshop objectives, scope and final architecture &
Complete \\
5.2 & Architecture & Detailed AWS architecture and request flows &
Complete \\
5.3 & Prerequisites & Required accounts, tools and local configuration &
Complete \\
5.4 & Source Code Preparation & Repository structure, environment
variables and validation & Complete \\
5.5 & AWS Infrastructure & Networking, IAM, EKS, ECR, ALB and CloudFront
& Complete \\
5.6 & Database Deployment & PostgreSQL, Redis and DynamoDB deployment &
Complete \\
5.7 & Backend Deployment & Backend, chat, dispatcher, AI service and
worker deployment & Complete \\
5.8 & Frontend Deployment & React build, S3 upload and CloudFront
distribution & Complete \\
5.9 & Monitoring and Logging & Logs, health checks, metrics and alarms &
Complete \\
5.10 & End-to-End Testing & Application, chat, AI and Lambda
verification & Complete \\
5.11 & Security and Cost & IAM, encryption, OIDC and cost assumptions &
Complete \\
5.12 & Troubleshooting & Resolved deployment and runtime failures &
Complete \\
5.13 & Clean-up & Safe resource deletion procedure & Complete \\
\end{longtable}
\endgroup
}

\vspace{0.5em}\noindent\textbf{Final outcome}\par

At the end of the workshop, the Internship Application Tracker is
available through CloudFront and supports:

\begin{itemize}
\tightlist
\item
  static frontend delivery
\item
  backend APIs
\item
  realtime chat
\item
  persistent data storage
\item
  asynchronous AI processing
\item
  transactional outbox publishing
\item
  Lambda-based event notifications
\item
  automated deployment through GitHub Actions
\end{itemize}
